% Fake It Til You Make It? -- arXiv manuscript scaffold.
%
% Structure and framing are fixed here BEFORE the experiments run, so the writing
% cannot quietly reorganise itself around whatever the numbers turn out to be.
% Every \RESULT{...} is a placeholder that must be filled from results/, and the
% macro renders in red so an unfilled one cannot survive a proofread.
\documentclass[11pt]{article}

\usepackage[margin=1in]{geometry}
\usepackage{amsmath,amssymb}
\usepackage{booktabs}
\usepackage{graphicx}
\usepackage{xcolor}
\usepackage[numbers,sort&compress]{natbib}
\usepackage{hyperref}
\hypersetup{colorlinks=true,linkcolor=blue,citecolor=blue,urlcolor=blue}

\newcommand{\RESULT}[1]{\textcolor{red}{\textbf{[#1]}}}
\newcommand{\TODO}[1]{\textcolor{orange}{\textit{[TODO: #1]}}}

\title{Fake It Til You Make It? Measuring the Real-to-Synthetic Exchange Rate
for Acne Severity Classification}

\author{%
  Gregory Savage\\
  \texttt{gregory.leonardo.savage@gmail.com}
}

\date{\today}

\begin{document}
\maketitle

\begin{abstract}
Diffusion-generated images are widely reported to improve medical image classifiers
when used as augmentation. Whether they can \emph{substitute} for real images, and at
what rate, is much less clear: the two questions are routinely conflated, and the
comparisons that would separate them---matched training budgets, a test set the
generator never saw, variance across seeds---are usually absent. We study acne
severity grading on ACNE04, training a fixed classifier on mixtures of real and
generated images at $0$, $25$, $50$, $75$ and $100$ percent synthetic while holding
the total number of training images constant, and evaluating every arm on the same
sealed real test set. We run two generator regimes that the literature tends to
merge: \emph{closed-set}, a diffusion model fine-tuned only on the real training
split, and \emph{open-set}, an off-the-shelf model that has never seen the dataset.
We report \RESULT{headline substitution result}, and estimate the exchange
rate---how many synthetic images are worth one real image---at \RESULT{k}. We further
find \RESULT{memorisation rate}, which bears directly on the claim that synthetic
dermatological faces are anonymous. Code, configurations and the test-set unsealing
audit log are released.
\end{abstract}

\section{Introduction}
\label{sec:intro}

Three questions are asked about generated training data, and only the first is
settled:

\begin{enumerate}
  \item \textbf{Augmentation.} Does adding synthetic images to real ones help?
        Broadly yes, including in dermatology \citep{sagers2022,ktena2024}.
  \item \textbf{Substitution.} Can synthetic images \emph{replace} real ones at a
        matched budget? Evidence from general vision says no, though the gap has
        narrowed \citep{sariyildiz2023,fan2024}.
  \item \textbf{Exchange rate.} How many synthetic images buy one real image, and
        does that number grow without bound? Barely measured \citep{wang2025}.
\end{enumerate}

\TODO{state the acne-specific gap: the single fully-synthetic acne prior reports
97.6\% training on synthetic and testing on real \citep{zein2024}, roughly 11
points above the strongest real-trained ACNE04 results we are aware of, without
budget control and with a synthetic healthy class. Frame replication under controls
as a contribution in its own right.}

\paragraph{Contributions.}
\begin{itemize}
  \item The first substitution curve for acne severity grading: performance as a
        continuous function of the synthetic fraction, at a matched training budget,
        against a sealed real test set.
  \item An estimate of the real-to-synthetic exchange rate $k(n)$ as a function of
        the real-data budget $n$, including the regime where $k$ is unbounded.
  \item A separation of closed-set from open-set generation, which we argue explains
        part of why published numbers in this area do not reconcile.
  \item Controls that bound the interpretation: a real/synthetic discriminability
        probe, a memorisation audit, a label-permutation control, and a real-subset
        control that sets the bar a synthetic arm must actually clear.
  \item A data-quality audit of ACNE04 itself (\S\ref{sec:audit}), independent of the
        synthetic-data question: 5.2\% byte-identical duplicates, 4.25\% mean
        train/test leakage in the dataset's own published folds, and 78.9\% grade
        agreement between the dataset's two annotations of the same photograph ---
        below several accuracies reported on it.
  \item Evidence that the ACNE04 severity label is a property of the annotation team
        rather than of the photograph: an independent expert re-annotation of 1{,}204 of
        the same images counts $3.2\times$ more lesions, ranks them differently
        ($\rho = 0.471$), and after being granted any monotone recalibration still
        agrees with the original grade on only 58\% of images.
\end{itemize}

\section{Related work}
\label{sec:related}
\TODO{condense docs/01\_related\_work.md. Three paragraphs: general vision; medical
and dermatology; acne specifically. Every number must be promoted to a full-text
verified citation first --- see docs/VERIFY.md.}

\section{The benchmark is not clean}
\label{sec:audit}

Every number reported on ACNE04 --- ours included --- is computed against labels and
splits that we found to be defective in ways no prior work we verified mentions. The
audit needs no model, no GPU and no network; it is reproducible in minutes from the
distributed archive with \texttt{scripts/audit\_acne04.py}.

\paragraph{Exact duplicates.}
Grouping all 1{,}457 files by MD5 of the file bytes yields 38 groups of size two:
76 files, 5.2\% of the corpus, are byte-identical copies under two filenames, usually
under two different \texttt{levleN} prefixes. Decoding to pixels and re-hashing
recovers exactly the same 38 groups.

\paragraph{The labels are derived, which turns the duplicates into a repeat-annotation
experiment.}
Each label row is \texttt{<file> <grade> <lesion\_count>}, and for all 1{,}457 records
the grade is exactly the Hayashi banding of the count (1--5, 6--20, 21--50, $>$50),
with no exceptions. The severity label therefore carries no information beyond the
count, and all label noise originates in counting. Because the duplicates were
annotated twice without the annotators knowing they were the same photograph, they
measure that noise directly:

\begin{center}
\begin{tabular}{lr}
\toprule
Exact lesion-count agreement & 12/38 \quad (31.6\%) \\
Median $|\Delta\text{count}|$ (max) & 1 \quad (16) \\
\textbf{Severity-grade agreement} & \textbf{30/38 \quad (78.9\%)} \\
95\% Wilson interval & 63.7--88.9\% \\
\bottomrule
\end{tabular}
\end{center}

All eight grade disagreements are cases where the two counts fall in different Hayashi
bands; four straddle the 5/6 boundary and four the 20/21 boundary. The latter are not
boundary jitter but differences of a factor of 2.4--3 (9 vs.\ 21, 8 vs.\ 24, 7 vs.\ 22,
9 vs.\ 24), and their filenames form consecutive runs, which is the signature of a
block of images ingested and counted twice by different procedures.

Published accuracies on this benchmark are 83.7\% \citep{wu2019}, 84.11\%
\citep{lds2024}, 84.52\% \citep{kieglfn2022} and 87.33\% \citep{ffpll2025}. They sit at or above the
dataset's own reproducibility on images it annotated twice. We are careful about what
this does and does not license: $n=38$ is small, the interval is wide, and duplicates
need not be a random sample of the corpus. The defensible reading is not that these
results are wrong, but that \emph{this benchmark cannot presently resolve differences
of a few points}, and that a reported improvement of 0.4 points is far inside the
label noise.

\paragraph{Two expert teams do not induce the same severity ordering.}
The duplicates measure reproducibility on 38 images. A far better powered measurement
is available. The AcneAI team \citep{acneai2024} independently re-annotated 1{,}204 of
the same photographs and released ACNE04-v2 with 32{,}443 lesion annotations against
the original's 18{,}983; original filenames are preserved, so the two annotations join
image by image.

They count \textbf{3.2$\times$ more lesions} (mean 26.9 vs.\ 8.4 per image; more on
90.1\% of shared images), and the share of images that ACNE04's own labelling function
would grade severe or worse moves from 11.7\% to 46.4\%. Pearson $r$ between the two
counts is 0.550 and Spearman $\rho$ is 0.471.

Part of this is definitional and entirely legitimate. Hayashi's bands are defined over
\emph{inflammatory} eruptions --- papules plus pustules --- and that paper reports
explicitly that global severity did not correlate with comedone counts
\citep{hayashi2008}. AcneAI annotate ``all lesions, either acne or non-acne lesions
that look like acne,'' including the tiny ones. \emph{The problem is not that the teams
disagree; it is that ACNE04 publishes a count and a grade without stating which lesions
were counted}, so the label depends on a convention the release never specifies.

The definitional account does not survive contact with the data. A constant change of
lesion definition would give a roughly constant ratio and preserve rank order. Instead
the per-image ratio has an interquartile range of 2.00--6.64, a 10th-to-90th-percentile
span of $11.5\times$, and falls monotonically with severity (median $8.00\times$ on
images v1 counted at 1--2 lesions, $2.00\times$ at 11--64); on 8.3\% of images v2
counted \emph{fewer}. To remove the definitional component entirely we grant the second
annotation an arbitrary monotone recalibration, re-banding its counts by quantile
matching so its grade histogram equals v1's by construction:

\begin{center}
\begin{tabular}{lcc}
\toprule
& grade agreement & QWK \\
\midrule
v2 counts under ACNE04's own Hayashi banding & 30.1\% & 0.296 \\
v2 counts after optimal monotone rescaling & \textbf{58.2\%} & \textbf{0.480} \\
\bottomrule
\end{tabular}
\end{center}

After the rescaling, 38.5\% of images differ by one grade and 3.2\% by two or more, and
$\rho$ does not recover on the images with the most lesions, where relative counting
noise should be smallest ($\rho = 0.376$ for v1 count $\geq 5$, $0.297$ for
$\geq 20$).

We are explicit about the construction: applying Hayashi bands to AcneAI's counts is
\emph{our} operation, not theirs --- they use a continuous 0--100 score and report
ICC 0.8 for their own method --- and it is clinically wrong on its face precisely
because their counts include comedones. That is the demonstration. It shows how much
the benchmark's published label depends on an unstated counting convention, and it does
not require deciding which team is right: two expert annotations of the same 1{,}204
photographs do not induce the same severity ordering, and ACNE04 gives no basis for
preferring one.

The consequence for the literature is not that reported accuracies are wrong but that
they measure something narrower than they are read as measuring: agreement with one
team's counting convention on one set of photographs. Differences of one to three
points between methods are differences in fitting that convention. This bears directly
on our own study, where it sets a noise floor below which a substitution effect is not
measurable (\S\ref{sec:results}).

\paragraph{The published splits leak.}
ACNE04 ships fixed five-fold splits (1{,}165 train / 292 test). Intersecting the
duplicate groups against them, \emph{every fold leaks}: on average 4.25\% of each test
split is byte-identical to an image in the corresponding training split --- 3.56\%
where the two labels agree, which a memorising model receives as free correct answers,
and 0.68\% where they conflict, which it cannot get right. Across all five folds, 62
distinct test images have a byte-identical twin in training. This is a systematic
upward bias of a few points inside every published number on these splits.

\paragraph{Filenames are not labels.}
42 of 1{,}457 files carry a \texttt{levleN} prefix that disagrees with their labelled
grade. Deriving labels from filenames --- the obvious shortcut given the naming scheme
--- corrupts 2.9\% of the dataset.

\paragraph{Deduplication instruments must be calibrated on this corpus, not inherited.}
Every ACNE04 image is a half-face shot at roughly $70^\circ$ against a dark background,
so generic perceptual and semantic embeddings have little signal to work with. At the
common perceptual-hash setting (Hamming $\leq 8$) linkage percolates into a single
cluster holding 14.7\% of the corpus, and inspection shows the linked pairs are
different people in the same pose; CLIP ViT-B/32 fails identically, with a median
pairwise cosine of 0.83 across the whole corpus and percolation beginning at 0.96.
We use Hamming $\leq 2$ and cosine $\geq 0.98$, which give 1{,}398 groups with a
largest cluster of three images, and we calibrate against the 38 byte-identical pairs,
all of which sit at cosine 1.0000.

\paragraph{What we do about it, and what remains open.}
We discard the published folds and split after deduplication on a frozen seed
(\S\ref{sec:method}), so our own test set is leak-free with respect to exact
duplication. Two consequences follow for reading our results. First, our real-only arm
should be expected to land \emph{below} the published band --- it has no leakage bonus,
it is scored with balanced accuracy on an imbalanced test set, and deduplication
removes training images --- and that is not evidence of a broken pipeline. Second,
78.9\% grade reproducibility is a noise floor: a substitution effect smaller than the
label noise is not an effect we can claim to have measured.

The audit is also incomplete in one specific way we flag rather than bury. Neither
perceptual hashing nor CLIP detects \emph{the same subject photographed twice}, and
ACNE04's capture protocol makes repeat subjects plausible. A face-identity embedding
would settle it. Until it is run, our splits and every published split carry an
unmeasured same-subject leakage risk.

\section{Method}
\label{sec:method}

\subsection{Task and data}
ACNE04 \citep{wu2019}: \RESULT{n images after deduplication}, Hayashi severity in
four ordinal grades. Deduplicated by perceptual hash and embedding cosine before
splitting, with whole duplicate clusters assigned to a single split
(\RESULT{n duplicate clusters}). Splits: \RESULT{n train} / \RESULT{n val} /
\RESULT{n test}.

\subsection{The sealed test set}
The real test split is used exactly once per reported model. It is excluded from
classifier training, early stopping, generator training, prompt selection, guidance
selection and hyperparameter search. This is enforced in code rather than by
convention, and every unsealing is appended to an audit log released with the paper
(\RESULT{n unsealings}).

\subsection{Generators}
\textbf{Closed-set.} \RESULT{base model}, LoRA rank \RESULT{rank}, fine-tuned on the
real training split only. \textbf{Open-set.} \RESULT{base model}, no fine-tuning,
prompted across age, sex, skin tone (Fitzpatrick I--VI), view, lighting and capture
style. Guidance scale selected on validation from
$\{1.5, 3, 5, 7.5, 10\}$: \RESULT{selected scale}.

\subsection{Mixtures}
For synthetic fraction $p$, the training set holds $(1-p)N$ real and $pN$ synthetic
images, with $N$ fixed and the real class histogram preserved. Real subsets are
nested across $p$. The additive design holds the real budget fixed and adds synthetic
at $k \in \{0,1,2,4,8,16\}$ times that budget, from which $k(n)$ is read off.

\subsection{Classifier}
ResNet-50, held fixed across arms. Initialisation is an experimental condition, not a
default: an ImageNet-pretrained backbone has consumed roughly 1.3M real photographs,
so a ``100\% synthetic'' arm built on one is not 100\% synthetic. Scratch and
pretrained results are reported separately and never pooled.

\section{Results}
\label{sec:results}

\subsection{Substitution at matched budget}
\begin{figure}[t]
  \centering
  % \includegraphics[width=0.75\textwidth]{../results/final/fig1_substitution.png}
  \caption{Balanced accuracy against synthetic fraction at a fixed training budget,
  for both generator regimes and both initialisations. Bands are 95\% intervals
  across \RESULT{n} seeds.}
  \label{fig:substitution}
\end{figure}

\RESULT{H1: 100\% synthetic vs 100\% real, with paired bootstrap CI}.
\RESULT{H2: trend slope and CI}.

\subsection{The exchange rate}
\begin{figure}[t]
  \centering
  % \includegraphics[width=\textwidth]{../results/final/fig2_additive.png}
  \caption{Additive scaling at four real budgets. $k(n)$ is the multiplier at which
  the synthetic curve reaches real-only performance at the next budget up.}
  \label{fig:additive}
\end{figure}

\RESULT{k(n) table, including budgets where k is unbounded}.

\subsection{The severe tail}
\RESULT{per-class recall for severe and very severe against synthetic fraction}.
Overall metrics can hold up while the rare grades collapse; the tail is where
mode-dropping shows first and where the clinical stakes are.

\subsection{Controls}
\textbf{Discriminability.} Real vs.\ synthetic AUC \RESULT{auc}, bounding how much of
any effect could be artefact-driven. \textbf{Memorisation.} \RESULT{rate}\% of
generated images exceed the replication threshold against the generator's training
set. \textbf{Real-subset control.} \RESULT{comparison} --- the bar a synthetic arm
must clear is real-only data of the same reduced size, not the full-data arm.
\textbf{Label permutation.} \RESULT{result}.

\subsection{Skin tone}
\RESULT{per-ITA-bin metrics}. Reported as a proxy: ITA measures pixels under a given
illuminant, not people, and acne erythema itself perturbs it.

\section{Discussion}
\label{sec:discussion}
\TODO{Write only after the numbers exist. If H1 holds, the contribution is the
exchange rate and the demonstration that the acne prior does not replicate under
budget control. If H1 is refuted, the contribution is larger and the controls in
\S\ref{sec:results} become the load-bearing part of the argument. Do not reorganise
the paper around whichever way it lands.}

\paragraph{The prototype effect.} A synthetic arm that \emph{beats} the real arm is
not by itself evidence that the generated images carry information the real data
lacked. A generator that reduces within-grade variation yields cleaner prototypes of
each class, and prototypes are easier to learn from. We encountered this directly:
a simulator that narrowed within-grade distributions while preserving the
label--image mapping produced synthetic-trained models that were measurably better.
It is a plausible mechanism for the 97.6\% reported by \citet{zein2024}. Any arm
exceeding the $p=0$ baseline is therefore checked against within-class embedding
variance, test-set difficulty stratification, and transfer to the external validation
set before being reported as a win. \RESULT{outcome of these checks}

\section{Limitations}
\label{sec:limitations}
\begin{itemize}
  \item Test set of \RESULT{n} images: powered for roughly 5-point differences in
        balanced accuracy, not 1--2 point differences.
  \item ACNE04 is a single-source corpus skewed towards lighter skin; the external
        validation set partly addresses this.
  \item Conclusions are tied to the generators evaluated. \citet{fan2024} attribute
        the synthetic-to-real gap to generator capability rather than to anything
        fundamental, so these numbers are a measurement at a point in time.
  \item ITA is a proxy for skin tone, not ground truth.
  \item No clinical claim is made. This is a methods paper about training data.
\end{itemize}

\section*{Ethics and data release}
ACNE04 is academic-use only and is not redistributed. Synthetic faces depicting a
visible dermatological condition are sensitive even when the subjects are fictional;
released images carry provenance metadata and are labelled synthetic. Release is
gated on the memorisation audit: if replication is material, we release metrics and
code but not images.

\section*{Reproducibility}
Code, configurations, split seeds, per-run provenance records and the test-set
unsealing audit log are released at \url{https://github.com/GLESAV/fake_it_til_you_make_it}.

\bibliographystyle{plainnat}
\bibliography{refs}

\appendix
\section{Guidance-scale sweep}
\RESULT{table: validation balanced accuracy by guidance scale, both regimes}.

\section{Prompt ablation}
\RESULT{minimal vs diversified prompts, open-set arm}.

\section{Architecture robustness}
\RESULT{ViT-B/16 and EfficientNet-B0 substitution curves}.

\end{document}
